\section{Related Work}

\subsection{LLM Role-Playing and Persona Simulation}

Recent work has demonstrated that LLMs can adopt diverse personas through prompting. Persona-driven generation explores how demographic and personality attributes affect LLM output. Our work extends this to the specific domain of academic expert review, where ``persona'' means a complete research identity with publications, institutional context, and domain-specific concerns.

\subsection{LLM-as-Judge}

The ``LLM-as-judge'' paradigm uses language models to evaluate outputs. Studies show that LLM judges exhibit biases (position bias, verbosity bias, self-enhancement bias). Our multi-reviewer panel approach mitigates individual biases through aggregation and explicitly measures inter-reviewer agreement to detect convergence.

\subsection{Multi-Agent Debate and Review}

Multi-agent systems for review and debate assign different roles to multiple LLM instances. Our contribution focuses on the \emph{calibration} question---not whether multiple agents can be created, but whether they produce genuinely distinct outputs that justify the multi-agent architecture.
